\section{Model Setup and Value Function Properties}

\subsection*{Question 1: Recursive Formulation of the American Option Value}

\textbf{Problem Statement:}
Prove that $V_t = v(t,S_t)$ such that:
\begin{align*}
v(T, x) &= h(x) \\
v(t-1, x) &= \max\left\{ h(x), \tilde{p} v(t, x(1+d)) + \tilde{q} v(t, x(1+u)) \right\}, \quad t=1,\dots,T
\end{align*}
where $\tilde{p}$ and $\tilde{q}$ are positive constants to be determined.

\textbf{Step-by-Step Solution:}
\begin{enumerate}
    \item \textbf{Terminal Condition:} At maturity $t=T$, the option expires. The value of the American option is simply its intrinsic payoff. Thus, $V_T = h(S_T)$, which implies $v(T, x) = h(x)$.
    
    \item \textbf{Dynamic Programming Principle:} At any time $t-1$, the holder of an American option must choose between two mutually exclusive actions:
    \begin{itemize}
        \item \textit{Exercise immediately:} The payoff received right now is $h(S_{t-1})$.
        \item \textit{Hold the option (Continuation):} The option is kept alive until at least time $t$. Its value is the discounted expected value of the option at time $t$ under the risk-neutral measure $Q$.
    \end{itemize}
    
    \item \textbf{Calculating the Continuation Value:} Because the market is complete and arbitrage-free, the Fundamental Theorem of Asset Pricing states that the expected discounted value under $Q$ is:
    \begin{equation*}
        \frac{1}{1+r} \mathbb{E}_Q[ V_t \mid \mathcal{F}_{t-1} ]
    \end{equation*}
    We know that $V_t = v(t, S_t)$. Under the binomial model, the asset price $S_t$ can transition from $S_{t-1}$ to:
    \begin{itemize}
        \item $S_{t-1}(1+u)$ with risk-neutral probability $q = \frac{r-d}{u-d}$
        \item $S_{t-1}(1+d)$ with risk-neutral probability $1-q$
    \end{itemize}
    Substituting these into the expectation gives the continuation value:
    \begin{equation*}
        \frac{1}{1+r} \left( (1-q) v(t, S_{t-1}(1+d)) + q v(t, S_{t-1}(1+u)) \right)
    \end{equation*}
    
    \item \textbf{Identifying the Constants:} Matching this expectation with the given formula yields:
    \begin{align*}
        \tilde{p} &= \frac{1-q}{1+r} \\
        \tilde{q} &= \frac{q}{1+r}
    \end{align*}
    Given the problem condition $-1 < d < r < u$, it follows that $0 < q < 1$ and $1+r > 0$. Therefore, both $\tilde{p}$ and $\tilde{q}$ are strictly positive.
    
    \item \textbf{Conclusion:} The rational option holder will choose the action that maximizes their payoff. Taking the maximum of the immediate exercise value and the continuation value yields the desired recursive relation:
    \begin{equation*}
        v(t-1,x) = \max\left\{ h(x), \tilde{p} v(t, x(1+d)) + \tilde{q} v(t, x(1+u)) \right\}
    \end{equation*}
\end{enumerate}


\subsection*{Question 2: Shape Preservation of the Value Function}

\textbf{Problem Statement:}
If $h$ is non-decreasing (resp. non-increasing, convex, continuous), then $x \mapsto v(t,x)$ is non-decreasing (resp. non-increasing, convex, continuous) for all $t=0,\dots,T$.

\textbf{Step-by-Step Solution:}
We prove this by backward induction on time $t$.

\begin{enumerate}
    \item \textbf{Base Case ($t=T$):} We have $v(T,x) = h(x)$. Since the payoff function $h(x)$ inherently possesses the property in question, $v(T,x)$ trivially holds the same property.
    
    \item \textbf{Inductive Hypothesis:} Assume that for some time $t \le T$, the function $x \mapsto v(t,x)$ possesses the required property (e.g., convexity). We must show that $x \mapsto v(t-1,x)$ also possesses it.
    
    \item \textbf{Analyzing the Continuation Function:} Let us define the continuation value function $C(x)$ at time $t-1$:
    \begin{equation*}
        C(x) = \tilde{p} v(t, x(1+d)) + \tilde{q} v(t, x(1+u))
    \end{equation*}
    \begin{itemize}
        \item From our base assumptions, $d > -1$, meaning $(1+d) > 0$ and $(1+u) > 0$. Scaling the input $x$ by a positive constant preserves monotonicity, convexity, and continuity. Thus, the terms $v(t, x(1+d))$ and $v(t, x(1+u))$ maintain the property.
        \item Since $\tilde{p} > 0$ and $\tilde{q} > 0$, taking a linear combination with strictly positive weights also preserves these properties. Therefore, $C(x)$ inherits the exact same shape properties as $v(t,\cdot)$.
    \end{itemize}
    
    \item \textbf{Applying Properties to the Maximum Operator:} 
    Recall the recursive relationship:
    \begin{equation*}
        v(t-1, x) = \max\{h(x), C(x)\}
    \end{equation*}
    A fundamental property of the max operator is that it preserves these specific functional shapes:
    \begin{itemize}
        \item The maximum of two non-decreasing functions is non-decreasing.
        \item The maximum of two non-increasing functions is non-increasing.
        \item The maximum of two convex functions is convex.
        \item The maximum of two continuous functions is continuous.
    \end{itemize}
    Since both the immediate payoff $h(x)$ and the continuation value $C(x)$ possess the required property, their maximum $v(t-1, x)$ does as well. This completes the inductive step, proving the statement for all $t=0,\dots,T$.
\end{enumerate}